\chapter{「スタック」タブの設定}\label{chap:stack}

スタックタブの設定は、\textbf{どのフレームを・どう重ねるか}を決めます。
ここの設定の多くは、変えても \ui{再スタック} だけで済みます（品質評価・アライメントはやり直しません）。
いろいろ試して比べるのに向いた設定です。

% ------------------------------------------------------------------------------
\section{スタック（画像の合成）}

\sidebyside[0.5]{insp_stack}{%
  \mk[e]{selmode}{1}\mk[e]{aptop}{2}\mk[e]{normalize}{3}\mk[e]{method}{4}\mk[w]{sigma}{5}\mk[e]{lowmem}{6}\mk[e]{rawcfa}{7}}{%
  \begin{marklist}
    \item フレームの選択方式
    \item 位置合わせ領域ごとに採用するフレーム
    \item 輝度正規化
    \item 加算方式
    \item σクリップの閾値
    \item 低メモリモード
    \item Bayer のまま合成する
  \end{marklist}
  \smallskip
  \uil{これだけを変えた再スタックは解析をやり直しません} とあるとおり、\badge{1}\,\badge{2} を変えても \ui{再スタック} だけで結果を見られます。}

\begin{setting}{\badge{1} フレームの選択方式}{割合}{割合／枚数}
  \badge{2} を「上位何\%」で決めるか、「上位何枚」で決めるかを選びます。
  \meyasu{動画の長さがまちまちでも同じ感覚で使えるのは \uil{割合}。ノイズの量をそろえたいとき（毎回200枚重ねる、など）は \uil{枚数}。}
\end{setting}

\begin{setting}{\badge{2} 位置合わせ領域ごとに採用するフレーム}{10 \%}{1〜100 \%（枚数のときは 1〜全フレーム数）}\label{sec:setting-aptop}
  位置合わせ領域ごとに、その場所の写りがよい順に何\%（何枚）を重ねるかです。\textbf{スタックの仕上がりをいちばん大きく左右する設定}です。
  \begin{center}
  \begin{tikzpicture}[font=\sffamily\scriptsize]
    \shade[left color=lsblue!70, right color=lsorange!80] (0,0) rectangle (9,0.35);
    \node[anchor=north west, align=left] at (0,-0.05) {少なくする（例 3〜5\%）\\よい瞬間だけ → シャープ\\重ねる枚数が減る → ざらつく};
    \node[anchor=north east, align=right] at (9,-0.05) {多くする（例 30〜50\%）\\悪い瞬間も混ざる → ぼやける\\枚数が増える → なめらか};
    \draw[white, line width=1.2pt] (1.8,0) -- (1.8,0.35); \node[above] at (1.8,0.35) {10\%（既定）};
  \end{tikzpicture}
  \end{center}
  \meyasu{重ねる枚数が\textbf{位置合わせ領域あたり 100〜500 枚}くらいになるのが1つの目安です（枚数＝使うフレーム数×割合。例：4617 フレームの 10\%＝約 460 枚）。
  \begin{itemize}[itemsep=0pt, topsep=1pt]
    \item 数千フレームの惑星動画：5〜15\%
    \item 数百フレームしかない、または暗くノイズが多い：20〜40\%
    \item ゆらぎがとても小さい日：多めにしてもシャープさが落ちにくい
  \end{itemize}
  迷ったら、「書き出し」欄の \uil{複数の採用率で書き出す} に \texttt{5, 10, 25} のように入れてまとめて作り、見比べるのが確実です（第\ref{sec:multiexport}節）。}
\end{setting}

\begin{setting}{\badge{3} 輝度正規化}{オン}{オン／オフ}
  薄い雲や空の透明度の変化で、フレームごとに明るさが変わることがあります。重ねる前に、各フレームの明るさをお手本にそろえます。
  \meyasu{オンのまま。明るさが本当に変わる現象を記録したいとき（ふつうは無い）だけオフにします。}
\end{setting}

\begin{setting}{\badge{4} 加算方式}{単純平均}{単純平均／品質重み付き平均／σクリップ}
  \begin{itemize}[itemsep=0pt, topsep=2pt]
    \item \uil{単純平均}：選んだフレームを同じ重みで平均します。
    \item \uil{品質重み付き平均}：品質が高いフレームほど重く数えます。シャープさを少し優先します。
    \item \uil{σクリップ}：画素ごとに、ほかのフレームと飛び抜けて違う値（人工衛星・飛行機の光跡、宇宙線、ホットピクセルなど）を除いて平均します。
  \end{itemize}
  \meyasu{ふだんは \uil{単純平均}。光跡や白い点が写り込んだときは \uil{σクリップ}（重ねる枚数が 20 枚以上あるときに効きます）。}
\end{setting}

\begin{setting}{\badge{5} σクリップの閾値（σ）}{2.0}{正の数（σクリップのときだけ使える）}
  平均からばらつきの何倍（σ）離れた値を除くかです。
  \meyasu{小さいほど多くの値を除きます。2.0〜3.0。光跡が消えきらないときは下げ、ノイズが増えたように見えるときは上げます。}
\end{setting}

\begin{setting}{\badge{6} 低メモリモード（2GB上限）}{オフ}{オン／オフ}
  スタック中に使うメモリをおよそ 2GB までに抑えます。結果は同じで、少し時間がかかることがあります。
  \meyasu{メモリの少ない Mac（4〜8GB）で、大きな画像やドリズル拡大のときに、ほかのアプリが重くなるならオンにします。}
\end{setting}

\begin{setting}{\badge{7} Bayer のまま合成する（デバイヤーしない）}{オフ}{オン／オフ}
  カラーカメラの画像を、色をそろえる前（Bayer 配列のまま）で重ね、そのまま書き出します。
  \meyasu{PixInsight など、ほかのソフトで後から色をそろえたい人向けです。ふだんはオフ。切り替えると品質評価からやり直しになります。}
\end{setting}

% ------------------------------------------------------------------------------
\section{ドリズル拡大}\label{sec:drizzle}

ドリズルは、フレームごとの細かな位置のずれを利用して、元の画素より細かい画素で画像を作り直す拡大の方法です。
\textbf{像が画素に対して粗く写っている（アンダーサンプリング）ときだけ}、解像度が上がります。
焦点距離が十分に長く、すでに細かく写っているときに倍率を上げても、ファイルが大きくなり処理が遅くなるだけです。

\sidebyside[0.5]{insp_drizzle}{\mk[e]{scale}{1}\mk[e]{pixfrac}{2}\mk[e]{diagnose}{3}\mkd[w]{result}{4}}{%
  \begin{marklist}
    \item 倍率
    \item 投影する画素の幅（pixfrac）
    \item \ui{ドリズルを診断}
    \item 診断の結果
  \end{marklist}
  \smallskip
  倍率を 1× 以外にすると、倍率の下に出力の大きさと作業に使うメモリの目安が出ます。}

\begin{setting}{\badge{1} 倍率}{1×（使わない）}{1×／1.5×／2×／3×}
  出力の画像を何倍の大きさで作るかです。変えたら \ui{再スタック}。
  \meyasu{まず \badge{3} の診断を押して、勧められた倍率にします。}
\end{setting}

\begin{setting}{\badge{2} 投影する画素の幅（pixfrac）}{0.90}{0.5〜1.0}
  1つの画素を、元の大きさの何割に縮めて新しい画素へ配るかです。小さいほどシャープになりますが、隙間ができてノイズやむらが出やすくなります。
  \meyasu{0.7〜0.9。重ねる枚数が多いほど小さくできます。まだらな模様が出たら大きくします。}
\end{setting}

\subsection*{\badge{3} ドリズルを診断}

アライメントのあとに \ui{ドリズルを診断} を押すと、スタックせずに数秒で「ドリズルが効きそうか・何倍まで効きそうか」の目安を出します。
次の3つを調べ、いちばん厳しいものを答えにします。

\begin{center}\small
\begin{tabularx}{0.95\linewidth}{@{}>{\sffamily\bfseries}l X@{}}
\toprule
像の細かさ & 画素で表せる最も細かい模様（ナイキスト周波数）を超えて、細部が写っているか。超えていれば、画素が粗すぎる＝ドリズルの出番です。\\
位置のばらつき & フレームごとの位置のずれが、1画素より細かい単位でばらけているか。ばらけていないと、細かい画素の隙間を埋められません。\\
採用枚数 & 位置合わせ領域あたりの枚数が足りるか。拡大すると1枚の寄与が細かい画素に分かれ、ノイズが増えます。\\
\bottomrule
\end{tabularx}
\end{center}

結果には、結論（\uil{効きそうです。目安は 2× までです。}または\uil{効果は小さそうです。等倍（1×）のままをおすすめします。}）と、
3つの理由が出ます。いちばん厳しかったものに\uil{（いちばんの制約）}と付きます。
効きそうなときは \ui{2× にする} のようなボタンが出て、押すと倍率を合わせます。
採用率を変えたり、アライメントをやり直したりすると、結果は消えます（もう一度押してください）。

\begin{caution}[診断は「ざっくりの目安」です]
  1枚の画像から見積もるので、正確な判定ではありません。確かめるには、倍率を変えてスタックし、拡大して見比べてください。
  背景の空が写っていない月面・太陽の全面では、控えめ（小さめ）に見積もります。
  上の画面の例では、木星の像はナイキストをわずかに超えるだけ（1.3倍）なので、1× を勧めています。
\end{caution}
